
    %%%%%%%%%%%%%%%%%%%%%%%%%%%%%%%%%%%%%%%%%%%%%%%%%%%%%%%%%%%%%%%%%%%%%%%%%%%%%%%%
    %%% Classe do documento
    \documentclass[12pt, addpoints]{exam}
    \UseRawInputEncoding

    %%%%%%%%%%%%%%%%%%%%%%%%%%%%%%%%%%%%%%%%%%%%%%%%%%%%%%%%%%%%%%%%%%%%%%%%%%%%%%%%
    % preambulo.tex
    %
    % Arquivo de configuração de packages para uso o LaTeX, conforme minhas
    % preferências e modelos pessoais.
    %
    % Para maiores informações, visite:
    %    https://github.com/abrantesasf/latex
    %
    % NÃO ALTERE SE NÃO SOUBER O QUE ESTÁ FAZENDO!
    %%%%%%%%%%%%%%%%%%%%%%%%%%%%%%%%%%%%%%%%%%%%%%%%%%%%%%%%%%%%%%%%%%%%%%%%%%%%%%%%


    %%%%%%%%%%%%%%%%%%%%%%%%%%%%%%%%%%%%%%%%%%%%%%%%%%%%%%%%%%%%%%%%%%%%%%%%%%%%%%%%
    %%% Estruturas de controle:
    \input{static/utils/controles}

    %%%%%%%%%%%%%%%%%%%%%%%%%%%%%%%%%%%%%%%%%%%%%%%%%%%%%%%%%%%%%%%%%%%%%%%%%%%%%%%%
    %%% Compilação condicional em PDF:
    \input{static/utils/pdf}

    %%%%%%%%%%%%%%%%%%%%%%%%%%%%%%%%%%%%%%%%%%%%%%%%%%%%%%%%%%%%%%%%%%%%%%%%%%%%%%%%
    %%% Configurações de layout da página:
    \input{static/utils/layout}

    %%%%%%%%%%%%%%%%%%%%%%%%%%%%%%%%%%%%%%%%%%%%%%%%%%%%%%%%%%%%%%%%%%%%%%%%%%%%%%%%
    %%% Configurações para autores:
    \input{static/utils/autores}

    %%%%%%%%%%%%%%%%%%%%%%%%%%%%%%%%%%%%%%%%%%%%%%%%%%%%%%%%%%%%%%%%%%%%%%%%%%%%%%%%
    %%% Cabeçalho e rodapé:
    %%% Atenção! É MUITO DIFÍCIL ajustar apropriadamente os running headers
    %%% e running footers da primeira vez. Você pode confiar no LaTeX e deixar que
    %%% ele faça o ajuste automaticamente ou pode quebrar a cabeça e
    %%% tentar melhorar algo que o LaTeX já faz muito bem, mesmo que não fique
    %%% exatamente do jeito que você quer. O que você prefere? Se quiser ajustar
    %%% manualmente, descomente a linha abaixo e faça as configurações no arquivo
    %%% "utils/headings.tex".
    %\input{static/utils/headings}

    %%%%%%%%%%%%%%%%%%%%%%%%%%%%%%%%%%%%%%%%%%%%%%%%%%%%%%%%%%%%%%%%%%%%%%%%%%%%%%%%
    %%% Configuração para as fontes do documento:
    %%% Se você quiser usar fontes próprias ou do sistema, precisará ajustar as
    %%% configurações no arquivo "utils/fontes.tex".
    %%% ATENÇÃO: por padrão eu utilizo XeLaTeX com fontes proprietárias projetadas
    %%% por Matthew Butterick, adquiridas em https://mbtype.com, que não podem ser
    %%% distribuídas devido à licença de utilização. Se você usar XeLaTeX, você
    %%% DEVERÁ OBRIGATORIAMENTE ajustar as fontes no arquivo "utils/fontes.tex".
    \input{static/utils/fontes}

    %%%%%%%%%%%%%%%%%%%%%%%%%%%%%%%%%%%%%%%%%%%%%%%%%%%%%%%%%%%%%%%%%%%%%%%%%%%%%%%%
    %%% Sumário:
    \input{static/utils/toc}

    %%%%%%%%%%%%%%%%%%%%%%%%%%%%%%%%%%%%%%%%%%%%%%%%%%%%%%%%%%%%%%%%%%%%%%%%%%%%%%%%
    %%% Matemática:
    \input{static/utils/matematica}

    %%%%%%%%%%%%%%%%%%%%%%%%%%%%%%%%%%%%%%%%%%%%%%%%%%%%%%%%%%%%%%%%%%%%%%%%%%%%%%%%
    %%% Notas de rodapé:
    \input{static/utils/footnotes}

    %%%%%%%%%%%%%%%%%%%%%%%%%%%%%%%%%%%%%%%%%%%%%%%%%%%%%%%%%%%%%%%%%%%%%%%%%%%%%%%%
    %%% Referências cruzadas, links, citações:
    \input{static/utils/crossrefs}

    %%%%%%%%%%%%%%%%%%%%%%%%%%%%%%%%%%%%%%%%%%%%%%%%%%%%%%%%%%%%%%%%%%%%%%%%%%%%%%%%
    %%% Configurações para os hiperlinks em PDF:
    \input{static/utils/pdfconfig}

    %%%%%%%%%%%%%%%%%%%%%%%%%%%%%%%%%%%%%%%%%%%%%%%%%%%%%%%%%%%%%%%%%%%%%%%%%%%%%%%%
    %%% Glossário e índice remissivo:
    \input{static/utils/glossindex}

    %%%%%%%%%%%%%%%%%%%%%%%%%%%%%%%%%%%%%%%%%%%%%%%%%%%%%%%%%%%%%%%%%%%%%%%%%%%%%%%%
    %%% Referências bibliográficas:
    \input{static/utils/bibliografia}

    %%%%%%%%%%%%%%%%%%%%%%%%%%%%%%%%%%%%%%%%%%%%%%%%%%%%%%%%%%%%%%%%%%%%%%%%%%%%%%%%
    %%% Cores:
    \input{static/utils/cores}

    %%%%%%%%%%%%%%%%%%%%%%%%%%%%%%%%%%%%%%%%%%%%%%%%%%%%%%%%%%%%%%%%%%%%%%%%%%%%%%%%
    %%% Computação:
    \input{static/utils/computacao}

    %%%%%%%%%%%%%%%%%%%%%%%%%%%%%%%%%%%%%%%%%%%%%%%%%%%%%%%%%%%%%%%%%%%%%%%%%%%%%%%%
    %%% Gráficos:
    \input{static/utils/graficos}

    %%%%%%%%%%%%%%%%%%%%%%%%%%%%%%%%%%%%%%%%%%%%%%%%%%%%%%%%%%%%%%%%%%%%%%%%%%%%%%%%
    %%% Ícones:
    \input{static/utils/icones}

    %%%%%%%%%%%%%%%%%%%%%%%%%%%%%%%%%%%%%%%%%%%%%%%%%%%%%%%%%%%%%%%%%%%%%%%%%%%%%%%%
    %%% Blocos:
    \input{static/utils/blocos}

    %%%%%%%%%%%%%%%%%%%%%%%%%%%%%%%%%%%%%%%%%%%%%%%%%%%%%%%%%%%%%%%%%%%%%%%%%%%%%%%%
    %%% Boxes coloridos:
    \input{static/utils/colorbox}

    %%%%%%%%%%%%%%%%%%%%%%%%%%%%%%%%%%%%%%%%%%%%%%%%%%%%%%%%%%%%%%%%%%%%%%%%%%%%%%%%
    %%% Tabelas:
    \input{static/utils/tabelas}

    %%%%%%%%%%%%%%%%%%%%%%%%%%%%%%%%%%%%%%%%%%%%%%%%%%%%%%%%%%%%%%%%%%%%%%%%%%%%%%%%
    %%% Ambientes de listas:
    \input{static/utils/listas}

    %%%%%%%%%%%%%%%%%%%%%%%%%%%%%%%%%%%%%%%%%%%%%%%%%%%%%%%%%%%%%%%%%%%%%%%%%%%%%%%%
    %%% Ambientes floats e similares:
    \input{static/utils/floats}

    %%%%%%%%%%%%%%%%%%%%%%%%%%%%%%%%%%%%%%%%%%%%%%%%%%%%%%%%%%%%%%%%%%%%%%%%%%%%%%%%
    %%% Ajustes para a classe EXAM:
    \input{static/utils/exam}

    %%%%%%%%%%%%%%%%%%%%%%%%%%%%%%%%%%%%%%%%%%%%%%%%%%%%%%%%%%%%%%%%%%%%%%%%%%%%%%%%
    %%% Ajustes para textos:
    \input{static/utils/texto}

    %%%%%%%%%%%%%%%%%%%%%%%%%%%%%%%%%%%%%%%%%%%%%%%%%%%%%%%%%%%%%%%%%%%%%%%%%%%%%%%%
    %%% Ambientes, aliases, comandos e customizações em geral para serem
    %%% utilizadas nos documentos. Defina aqui qualquer customização específica
    %%% para seus documentos!
    \input{static/utils/customizacoes}

    %%%%%%%%%%%%%%%%%%%%%%%%%%%%%%%%%%%%%%%%%%%%%%%%%%%%%%%%%%%%%%%%%%%%%%%%%%%%%%%%
    %%% Ajuste do layout, espaçamento de linhas e etc.:
    \geometry{a4paper, portrait, left=2.0cm, right=2.0cm, top=2.0cm, bottom=2.0cm}
    \singlespacing

    %%%%%%%%%%%%%%%%%%%%%%%%%%%%%%%%%%%%%%%%%%%%%%%%%%%%%%%%%%%%%%%%%%%%%%%%%%%%%%%%
    %%% Configurações de cabeçalho e rodapé (não use fancyhdr pois ocorre conflito):
    \pagestyle{headandfoot}
    \runningheadrule
    \firstpageheader{}{}{}
    \runningheader{Sem disciplina}{}{Avaliação}
    \firstpagefooter{}{}{Página \thepage\ de \numpages}
    \runningfooter{}{}{Página \thepage\ de \numpages}
    \runningfootrule

    %%%%%%%%%%%%%%%%%%%%%%%%%%%%%%%%%%%%%%%%%%%%%%%%%%%%%%%%%%%%%%%%%%%%%%%%%%%%%%%%
    %%% Configurações para as propriedades do PDF:
    \hypersetup{
    hidelinks,           % Comente para web, descomente para impressão
    colorlinks=true,      % True para web, False para impressão
    pdftitle=Ai: Avaliação,
    pdfauthor=Matheus Endlich Silveira,
    pdfsubject=Sem disciplina,
    pdfkeywords=['Sem tag'],
    pdfinfo={
        CreationDate={}, % Ex.: D:AAAAMMDDHH24MISS
        ModDate={}       % Ex.: D:AAAAMMDDHH24MISS
    }
    }
    \input{static/utils/pdfconfig}

    %%%%%%%%%%%%%%%%%%%%%%%%%%%%%%%%%%%%%%%%%%%%%%%%%%%%%%%%%%%%%%%%%%%%%%%%%%%%%%%%
    %%% Começa o documento
    \begin{document}

    %%%%%%%%%%%%%%%%%%%%%%%%%%%%%%%%%%%%%%%%%%%%%%%%%%%%%%%%%%%%%%%%%%%%%%%%%%%%%%%%
    %%% Folha de instruções padronizadas da UVV para a prova
    \pdfbookmark[1]{Instruções}{instrucoes}
    \begin{figure}[H]
        \begin{center}
            \includegraphics[scale=0.3]{{static/imgs/logo-uvv.png}}
        \end{center}	
    \end{figure}

    \vspace{-1cm}

    \begin{table}[H]
        \centering
        \begin{tabular}{|llll|}
            \hline
            \multicolumn{2}{|l|}{\textbf{Disciplina:} Sem disciplina} & 
            \multicolumn{1}{l|}{\multirow{3}{*}{\textbf{\makecell{Nota:\\ \,\\ \,}}}} & 
            \multirow{3}{*}{\textbf{\makecell{Coordenador:\\ \,\\ \,}}} \\ 
            \cline{1-2}
            \multicolumn{2}{|l|}{\textbf{Professor:} Matheus Endlich Silveira \hspace{4.72cm} } & 
            \multicolumn{1}{l|}{} & \\ 
            \cline{1-2}
            \multicolumn{2}{|l|}{\textbf{Aluno:}} & 
            \multicolumn{1}{l|}{} & \\ 
            \hline
            \multicolumn{1}{|l|}{\textbf{Turma:} \hspace{6.3cm} } & 
            \multicolumn{1}{l|}{\textbf{Semestre:}} & 
            \multicolumn{2}{l|}{\textbf{Valor:} 10 pontos} \\ 
            \hline
            \multicolumn{1}{|l|}{\textbf{Data:}} & 
            \multicolumn{3}{l|}{\textbf{Avaliação:}} \\ 
            \hline
        \end{tabular}
    \end{table}

    \begin{center}
        \fbox{\setlength\fboxsep{0.3cm} \fbox{\parbox{15.4cm}{
            \begin{center}
                \textbf{INSTRUÇÕES PARA A REALIZAÇÃO DESTA AVALIAÇÃO:}
            \end{center}
            \begin{itemize}[leftmargin=*]
                \item Esta é a primeira avaliação da disciplina Sem disciplina, 
                    e engloba o conteúdo referente Sem tag.
                \item Esta avaliação contém 20 questões objetivas, \textbf{todas obrigatórias}.
                \item \textbf{Leia atentamente cada questão!} As questões objetivas terão uma,
                    e apenas uma única, resposta a ser marcada.
                \item \textbf{Desligue o celular} antes de começar. A avaliação é
                    \textbf{individual} e \textbf{sem consulta}.
                \item As questões podem ser respondidas, na prova, com lápis ou caneta. Ao final
                    da prova \textbf{{VOCÊ DEVERÁ PREENCHER O GABARITO COM CANETA
                    DE COR PRETA}} \textbf{\underline{OU AZUL}} para que seja feita a correção
                    automatizada através da leitura do padrão de respostas marcado no
                    gabarito.
                \item \textbf{CUIDADO ao preencher o gabarito} pois eles são \textbf{INDIVIDUAIS
                    e NOMEADOS} (cada estudante tem um gabarito próprio específico, com um
                    código de identificação único). \textbf{Questões rasuradas são anuladas}
                    pelo software de leitura do gabarito.
                \item Ao final da prova, \textbf{DEVOLVA A PROVA E O GABARITO} para o professor.
                \item Siga todas as normas de \textbf{Integridade Acadêmica} da disciplina.
                    Alunos que forem flagrados com qualquer espécie de ``cola'' ou trocando
                    informações com outros alunos terão suas avaliações recolhidas, as notas
                    zeradas, e a situação será encaminhada para a coordenação para a aplicação
                    das penalidades previstas pela universidade.
                \item Com 90 minutos de avaliação você terá tempo de até 4,min por questão,
                    em média.
                \item Boa avaliação!
            \end{itemize}
            \vspace{2cm}
        }}}
    \end{center}
    \newpage

    %%%%%%%%%%%%%%%%%%%%%%%%%%%%%%%%%%%%%%%%%%%%%%%%%%%%%%%%%%%%%%%%%%%%%%%%%%%%%%%%
    %%% Inicia o bloco de questões:
    \begin{questions}

    \newpage
    \question

No que diz respeito as vantagens da arquitetura de micro-núcleo para sistemas operacionais em relação a arquiteturas de núcleo monolítico, quais das seguintes afirmações são verdadeiras?



\begin{itemize}

    \item[I] A arquitetura de micro-núcleo facilita a depuração do SO.

    \item[II] A arquitetura de micro-núcleo permite um número menor de mudanças de contexto.

    \item[III] A arquitetura de micro-núcleo facilita a reconfiguração de serviços do SO pois a maioria deles reside em espaço de usuário.

\end{itemize}
\begin{choices}
\choice Todas são verdadeiras
\choice I e III
\choice Apenas I
\choice I e II
\choice II e III
\end{choices}

\question

Dentre as características do modelo relacional e do modelo de objetos em bancos de dados, qual afirmação é \textbf{INCORRETA}?


\begin{choices}
\choice O modelo de objetos provê uma representação bem próxima de linguagens de programação.
\choice O modelo de objetos é mais adequado para a representação de tipos abstratos de dados.
\choice O modelo de objetos possui mais recursos estruturais para a representação de dados que o relacional.
\choice O relacionamento de herança é diretamente representado no modelo relacional.
\choice O relacionamento binário \( N \times M \) é representado de modo semelhante nos dois modelos.	
\end{choices}

\question

São linguagens clássicas que representam aplicações científicas, aplicações 

comerciais, aplicações de inteligência artificial, e aplicações gerais,

respectivamente:
\begin{choices}
\choice R, COBOL, C, Python
\choice COBOL, Java, LISP, Fortran
\choice Fortran, COBOL, LISP, C
\choice MATLAB, COBOL, SML, LISP
\choice Java, MATLAB, Scheme, JavaScript
\end{choices}

\question

Numa prova de múltipla escolha com 10 questões e 4 alternativas, qual a chance (probabilidade) de um aluno apenas "chutando as respostas" conseguir "gabaritar" a prova (acertar todas as questões).


\begin{choices}
\choice \( \frac{1}{2^{20}} \)
\choice \( \frac{1}{4^{15}} \)
\choice \( \frac{1}{10^{8}} \)
\choice \( \frac{1}{4^{20}} \)
\choice \( \frac{1}{10^{4}} \)
\end{choices}

\question

Uma prova de vestibular foi elaborada com 25 questões de múltipla escolha com 5 alternativas. O número de candidatos presentes à prova foi 63127. Considere a afirmação: Pelo menos 2 candidatos responderam de modo idêntico as \( k \) primeiras questões da prova. Qual é o maior valor de \( k \) para o qual podemos garantir que a afirmação é verdadeira.


\begin{choices}
\choice 10
\choice 8
\choice 9
\choice 6
\choice 7
\end{choices}

\question

Considere as seguintes tabelas em uma base de dados relacional (chaves primárias sublinhadas):\\

Departamento (\underline{CodDepto}, NomeDepto)\\

Empregado (\underline{CodEmp}, NomeEmp, CodDepto)\\

Considere as seguintes restrições de integridade sobre esta base de dados relacional:\\

- Empregado.CodDepto é sempre diferente de NULL\\

- Empregado.CodDepto é chave estrangeira da tabela Departamento com cláusulas ON DELETE RESTRICT e ON UPDATE RESTRICT\\

Qual das seguintes validações não é especificada por estas restrições de integridade:
\begin{choices}
\choice Sempre que uma nova linha for inserida em Departamento, deve ser garantido que o valor de Departamento.CodDepto aparece na coluna Empregado.CodDepto.
\choice Sempre que uma linha for excluída de Departamento, deve ser garantido que o valor de Departamento.CodDepto não aparece na coluna Empregado.CodDepto.
\choice Sempre que o valor de Empregado.CodDepto for alterado, deve ser garantido que o novo valor de Empregado.CodDepto aparece em Departamento.CodDepto.
\choice Sempre que o valor de Departamento.CodDepto for alterado, deve ser garantido que não há uma linha com o antigo valor de Departamento.CodDepto na coluna Empregado.CodDepto.
\choice Sempre que uma nova linha for inserida em Empregado, deve ser garantido que o valor de Empregado.CodDepto aparece na coluna Departamento.CodDepto.
\end{choices}

\question

O determinante da matriz dada abaixo é

$\begin{pmatrix}

2 & 7 & 9 & -1 & 1 \\

2 & 8 & 3 & 1 & 0 \\

-1 & 0 & 4 & 3 & 0 \\

2 & 0 & -1 & 0 & 0 \\

3 & 0 & 0 & 0 & 0 \\

\end{pmatrix}$
\begin{choices}
\choice -96
\choice -86
\choice 96
\choice 86
\choice 46
\end{choices}

\question

Considere as afirmações abaixo sobre endereços IP:

\begin{itemize}

    \item[I.] Uma rede IP classe C fornece até 256 endereços válidos para serem atribuídos a equipe.

    \item[II.] A quantidade máxima de bits que pode ser utilizada para se definir sub-redes em uma rede IP classe C é seis (6).

    \item[III.] A máscara padrão para uma rede classe B é 255.255.255.0.

\end{itemize}

Qual das seguintes alternativas representa as afirmações corretas?
\begin{choices}
\choice Somente II e III.
\choice Somente I.
\choice Somente III.
\choice Somente I e II.
\choice Somente II.
\end{choices}

\question

O contador da figura abaixo é:

\begin{figure}

    \centering

    \includegraphics[width=0.5\linewidth]{static/imgs/defaul.png}

    \label{fig:enter-label}

\end{figure}
\begin{choices}
\choice assíncrono
\choice isócrono
\choice síncrono
\choice anisócrono
\choice auto-sincronizado
\end{choices}

\question

Qual das seguintes condições não é necessária para a ocorrência de um \textit{deadlock}?


\begin{choices}
\choice Escalonamento preemptivo de recursos.
\choice Processos em espera circular.
\choice Uso mutuamente exclusivo de recursos por processos.
\choice Alocação parcial de recursos a processos.
\choice Haver compartilhamento de recursos por processos.
\end{choices}

\question

Sobre a técnica conhecida como Z-buffer é correto afirmar que:
\begin{choices}
\choice É uma técnica muito comum de detecção de colisão. 
\choice É possível realizar o cômputo das variáveis envolvidas de forma incremental.
\choice Nenhuma das alternativas acima está correta. 
\choice As dimensões do Z-buffer são independentes das dimensões do frame buffer. 
\choice As primitivas geométricas precisam estar ordenadas de acordo com a distância em relação ao observador. 
\end{choices}

\question

Considere o seguinte código para implementar exclusão mútua entre dois processos \(i\) e \(j\): 

Processo \(P_i\)

\begin{lstlisting}[language=C, basicstyle=\ttfamily\small]

do {

    while (turn != i); 

    // secao critica

    secao critica

    turn = j;

    // saída da secao critica

    codigo restante

} while (1);

\end{lstlisting} 
\begin{choices}
\choice Exige alternância estrita. 
\choice Os processos fazem espera ativa.
\choice A implementação garante exclusão mútua.
\choice Um processo bloqueia o outro mesmo não estando na seção crítica.
\choice A implementação garante progresso. 
\end{choices}

\question

Na definição da aquisição de um novo software de

gerenciamento de bancos de dados (SGBD) para uma empresa da área de transporte 

coletivo urbano, a direção da área de informática conduziu o processo de decisão

da seguinte forma: foi designado um profissional da área de bancos de dados

(aquele com maior experiência na área) e atribuída a ele a tarefa de decidir

qual seria o melhor SBGD a ser adquirido. Esse profissional desenvolveu uma

série de estudos sobre as opções disponíveis, utilizando ténicas de simulação e

testes específicos, levando em conta que a quase totalidade dos dados que serão

armazenados são dados de cadastro altamente estruturados e que devem estar

disponíveis na Internet para consulta pelos usuários. Ao final, apresentou

ao diretor um relatório em que indicava claramente qual o melhor tipo de SGBD a

ser adquirido pela empresa. Com base nessa informação, o diretor da empresa

disparou o processo de compra do SGBD. Podemos concluir que o profissional de

bancos de dados escolheu um:
\begin{choices}
\choice SGBD Hierárquico, para armazenar corretamente e de forma otimizada todas os relacionamentos entre os ônibus, empregados e passageiros do transporte público.
\choice SGBD Grafo, para armazenar todos os relacionamentos entre os usuários do transporte público nos ônibus.
\choice SGBD Relacional, para armazenar de forma otimizada os dados que, em sua maioria, são altamente estruturados.
\choice Nenhuma das respostas acima.
\choice SGBD Não Relacional (NoSQL) já que, para armazenar todos os dados e ainda permitir a consulta pelos usuários na Internet, o melhor seria utilizar uma solução com bancos de dados não relacionais que são otimizados para uso na web.
\end{choices}

\question

Uma integração de Sistemas Computacionais formando uma rede, tipicamente é implementada através da instalação de uma Arquitetura de Rede, que é composta de camadas e protocolos, em cada um dos elementos que compõem esta rede. Considere que estações “conversam” quando aplicações de usuários conseguem comunicar-se, sintática e semanticamente, através da Rede de Computadores. Baseados nesta premissa e em todos os conceitos associados à implementação e utilização das redes de computadores podemos afirmar como certo:


\begin{choices}
\choice Computadores com arquiteturas diferentes podem “conversar” através de multiplexadores.
\choice Computadores com arquiteturas de redes diferentes podem “conversar” através de um \textit{gateway} ou conversor de protocolos.
\choice Computadores com arquiteturas de redes diferentes conseguem “conversar”.
\choice Nenhuma delas é uma afirmação correta.
\choice Computadores com arquiteturas de rede parecidas conseguem “conversar”.
\end{choices}

\question

Qual dos protocolos abaixo pode ser caracterizado como protocolo de roteamento do tipo estado de enlace?
\begin{choices}
\choice ICMP
\choice RIP2
\choice OSPF
\choice IGMP
\choice BGP-4
\end{choices}

\question

Sejam os seguintes predicados de uma linguagem de primeira ordem:



\begin{itemize}

    \item \( N(x) \): \( x \) é número;

    \item \( P(x) \): \( x \) tem propriedade \( P \);

    \item \( x < y \): \( x \) é menor que \( y \).

\end{itemize}



E sejam os símbolos:

\begin{itemize}

    \item \( \forall \): quantificador universal;

    \item \( \Rightarrow \): operador se-então;

    \item \( \neg \): operador de negação.

\end{itemize}



Para a fórmula: \( \forall x (N(x) \Rightarrow \neg \forall y (N(y) \Rightarrow y < x)) \), qual alternativa abaixo \textbf{NÃO} constitui uma tradução possível?


\begin{choices}
\choice Para todo número, existe um outro número que é maior do que ele.
\choice Para qualquer \( x \), se \( x \) é número, então não é verdade que todos os números são menores do que ele.
\choice Não há um número tal que todos os números são menores do que ele.
\choice Para todo número, não é verdade que qualquer número seja menor do que ele.
\choice Não há um número menor do que outro número.
\end{choices}

\question

A velocidade de um ponto em movimento é dada pela equação\\

$v(t) = t e^{-0.01t} \, \text{m/s}$\\

O espaço percorrido desde o instante que o ponto começou a se mover até a sua parada total é


\begin{choices}
\choice \( 10^2 e^{-1} \, m \)
\choice \( 10^4 \, m \)
\choice \( 10^3 e^{-0.01} \, m \)
\choice \( 10^2 \, m \)
\choice \( (e^{-100} - 1) m \)
\end{choices}

\question

Seu chefe que fazer uma campanha de venda para produtos que custam entre 20 e 30

reais, mas está preocupado porque ouviu de um dos vendedores que diversos

produtos nessa faixa de preço estão com o estoque zerado. Para verificar a

situação seu chefe pediu para que você prepare um relatório que mostre os

produtos que custam entre 20 e 30 reais e que estão com o estoque zerado. Você

preparou o seguinte relatório para seu chefe:

\begin{tcolorbox}

 \begin{verbatim}1  SELECT l.nome AS loja

2       , p.nome AS produto

3       , p.preco_unitario AS preco

4       , e.quantidade AS qtd

5  FROM produtos p

6  INNER JOIN estoques e ON (e.produto_id = p.produto_id)

7  INNER JOIN lojas l ON (e.loja_id = l.loja_id)

8  WHERE (p.preco_unitario >= 20 OR p.preco_unitario <= 30)

9    AND e.quantidade = 0

10 ORDER BY loja, produto;

 \end{verbatim}

\end{tcolorbox}

O seu relatório:
\begin{choices}
\choice Está correto e trará exatamente, os produtos com o estoque zerado e na faixa de preço requisitada pelo seu chefe. Além disso ordena o resultado pelo nome da loja e pelo nome do produto.
\choice Está errado pois a tabela que deveria estar na cláusula FROM é a tabela ``lojas', não a tabela ``produtos'.
\choice Está errado pois há um erro de sintaxe nas linhas 6 e 7.
\choice Nenhuma das alternativas anteriores.
\choice Está errado pois vai trazer também produtos fora da faixa de preço que o seu chefe quer.
\end{choices}

\question

O número de \textit{strings binárias} de comprimento 7 e contendo um par de zeros consecutivos é


\begin{choices}
\choice 94
\choice 92
\choice 90
\choice 91
\choice 95
\end{choices}

\question

São motivos importantes e gerais para o estudo das linguagens de programação,

exceto:
\begin{choices}
\choice Aumentar a capacidade de expressar idéias
\choice Aprender a sintaxe e a semântica de uma determinada linguagem
\choice Melhorar o embasamento na escolha das linguagens para um projeto
\choice Avanço geral da computação
\choice Entender a importância da implementação
\end{choices}



    %%%%%%%%%%%%%%%%%%%%%%%%%%%%%%%%%%%%%%%%%%%%%%%%%%%%%%%%%%%%%%%%%%%%%%%%%%%%%%%%
    %%% Finaliza o bloco de questões:
    \end{questions}
    \end{document}
    